In this thesis, we have developed a unified mathematical and computational framework for evaluating high-dimensional optimal transport under structural constraints. Chapter~\ref{ch:dykstra} presents a theoretical resolution to the stalling phenomenon in Dykstra's algorithm, introducing the fast-forward algorithmic modification. Concurrently, Chapter~\ref{ch:optimal_transport} introduces a density tracking methodology that learns KR triangular maps using a probabilist's Hermite polynomial basis, which is then integrated into the scalable inexact projection framework detailed in Chapter~\ref{ch:optimisation}. In this chapter, we demonstrate the application of the theoretical framework to several examples to illustrate the effectiveness of our approach.

The numerical experiments are designed to validate the representational capacity of the transport maps and to assess the computational scalability of the custom optimisation solver. The initial evaluations examine the framework's geometric reconstruction capabilities using non-linear synthetic distributions. Subsequently, the analysis focuses on the primary aerospace application, specifically the transformation of cislunar satellite uncertainty parameterised by equinoctial elements. To further assess the methodology in the context of chaotic dynamical systems relevant to geophysics, the framework is evaluated using a dataset emanating from the Lorenz-63 model. Finally, the analysis returns to the core optimisation architecture to quantify the runtime efficiency of the inexact projected gradient descent solver in comparison with several existing solvers. The results of these experiments are reproducible using the second of two repositories accompanying this thesis~\cite{DykstraProjectRepo, DykstraOptimalTransportRepo}, available at:

\begin{center}
    \href{https://github.com/ClouD-161803/Optimal-Transport-Dykstra}{ClouD-161803/Optimal-Transport-Dykstra.git}.
\end{center}

\section*{Numerical setup}

For each of the experiments described in this chapter, the numerical setup implements the architecture outlined in Chapters~\ref{ch:dykstra}--\ref{ch:optimal_transport}--\ref{ch:optimisation}. To initialise the experiment, the state dimension $D$, number of particles $n$, and random seed are selected. The weights $w^{(0)}$ are initialised to yield the identity map $S^k(x^{(i)}) = x^{(i)}$. A maximum Hermite polynomial degree $J_k + 1$ is selected. Matrix $A$ and monotonicity tolerance vector $\epsilon_m$ are initialised according to~\eqref{eq:discrete_monotonicity}. The outer projected gradient descent loop executes for $T$ total iterations to update the decision variables $w^{(t)}$. At each outer loop iteration, the objective is computed according to~\eqref{eq:finite_objective}. To manage the computational demand of evaluating the full set of $n$ particles, stochastic mini-batching is utilised with a subset cardinality $B$. The learning rate dictates the magnitude of the unconstrained step to $w^\circ$ and is updated at each iteration, starting from an initial value $\eta^\circ$ and decaying according to the parameter $\gamma$. Iterative hard thresholding is applied periodically at every $K_{\mathrm{IHT}}$-th iteration, where non-physical weights are pruned via the boolean active mask vector $\mathcal{M}^{(t)}$ according to the threshold $\epsilon_p$. Within the inner solver cycle, the Dykstra projection budget ceiling $M^{(t)}$ dynamically scales based on the base budget $M_{\mathrm{base}}$ and the scaling exponent $p$. Given the full set of initialisations, we run the combined routine Algorithm~\ref{alg:fastforward}--Algorithm~\ref{alg:inexact_pgd}--Algorithm~\ref{alg:accelerated_dykstra}.